%% =============================================================
%%  ASCLab @ RPI — Technical / Research Report Template
%%  Advanced Space Concepts Laboratory
%%  Rensselaer Polytechnic Institute
%%


\documentclass[11pt, letterpaper]{article}

%% ---- PACKAGES -----------------------------------------------
\usepackage[top=1in, bottom=1in, left=1in, right=1in]{geometry}
\usepackage{graphicx}
\usepackage{amsmath, amssymb, amsthm}
\usepackage{booktabs}           % professional tables
\usepackage{hyperref}
\usepackage{xcolor}
\usepackage{fancyhdr}
\usepackage{titlesec}
\usepackage{caption}
\usepackage{subcaption}
\usepackage{natbib}             % AAS/AIAA-style references; use \citep, \citet
\usepackage{setspace}
\usepackage{microtype}
\usepackage{enumitem}
\usepackage{multirow}
\usepackage{array}
\usepackage{float}
\usepackage[T1]{fontenc}
\usepackage{lmodern}
\usepackage{lastpage}           % for "Page X of Y"

%% ---- COLORS (RPI brand) ------------------------------------
\definecolor{RPIred}{HTML}{D6001C}
\definecolor{RPIdark}{HTML}{222222}
\definecolor{ASCblue}{HTML}{1B3A6B}   % ASCLab accent navy

%% ---- HYPERLINK STYLE ----------------------------------------
\hypersetup{
    colorlinks = true,
    linkcolor  = ASCblue,
    citecolor  = ASCblue,
    urlcolor   = RPIred,
}

%% ---- SECTION FORMATTING -------------------------------------
\titleformat{\section}
    {\color{ASCblue}\large\bfseries}{\thesection.}{0.6em}{}[\color{ASCblue}\titlerule]
\titleformat{\subsection}
    {\color{RPIdark}\normalsize\bfseries}{\thesubsection.}{0.5em}{}
\titleformat{\subsubsection}
    {\color{RPIdark}\normalsize\itshape}{\thesubsubsection.}{0.5em}{}

%% ---- THEOREM ENVIRONMENTS -----------------------------------
\newtheorem{theorem}{Theorem}[section]
\newtheorem{lemma}[theorem]{Lemma}
\newtheorem{proposition}[theorem]{Proposition}
\newtheorem{corollary}[theorem]{Corollary}
\newtheorem{definition}{Definition}[section]
\newtheorem{remark}{Remark}[section]

%% ============================================================
%%  DOCUMENT METADATA — FILL IN BELOW
%% ============================================================
\newcommand{\ReportTitle}{Title of Technical Report}
\newcommand{\ReportSubtitle}{Subtitle or Project Name}        % leave blank if none: {}
\newcommand{\ReportNumber}{ASCLab-TR-2026-XX}                 % internal report number
\newcommand{\AuthorList}{Author One, Author Two, Author Three}
\newcommand{\AffiliationText}{%
    Advanced Space Concepts Laboratory (ASCLab)\\
    Department of Mechanical, Aerospace, and Nuclear Engineering\\
    Rensselaer Polytechnic Institute, Troy, NY 12180
}
\newcommand{\CorrespondingEmail}{placeholder@rpi.edu}
\newcommand{\ReportDate}{\today}
\newcommand{\Keywords}{cislunar SDA; trajectory optimization; space domain awareness; CR3BP}
\newcommand{\SponsorText}{%
    This work was supported in part by the xxxx.                % update or remove as appropriate
}
%% ============================================================

%% ---- HEADER / FOOTER ----------------------------------------
\pagestyle{fancy}
\fancyhf{}
\fancyhead[L]{\small\textcolor{ASCblue}{\textbf{ASCLab} @ RPI}}
\fancyhead[R]{\small\textcolor{RPIdark}{\ReportNumber}}
\fancyfoot[C]{\small\textcolor{RPIdark}{Page \thepage\ of \pageref{LastPage}}}
\fancyfoot[R]{\small\textcolor{RPIdark}{\ReportDate}}
\renewcommand{\headrulewidth}{0.4pt}
\renewcommand{\headrule}{\hbox to\headwidth{\color{RPIred}\leaders\hrule height \headrulewidth\hfill}}
\renewcommand{\footrulewidth}{0.2pt}

%% =============================================================
\begin{document}
\onehalfspacing

%% ---- COVER PAGE ---------------------------------------------
\begin{titlepage}
\newgeometry{top=0.75in, bottom=0.75in, left=1in, right=1in}
\vspace{0.3cm}
{\color{RPIred}\rule{\linewidth}{2pt}}
\vspace{0.5cm}

% -- Report number
\begin{center}
{\large\color{RPIdark}\texttt{\ReportNumber}}\\[0.5cm]

% -- Title
{\Huge\bfseries\color{ASCblue}\ReportTitle}\\[0.4cm]
\ifx\ReportSubtitle\empty\else
    {\Large\color{RPIdark}\ReportSubtitle}\\
\fi
\vspace{1.2cm}

% -- Authors
{\large \AuthorList}\\[0.4cm]
{\normalsize\color{RPIdark} \AffiliationText}\\[0.3cm]
{\normalsize Correspondence: \href{mailto:\CorrespondingEmail}{\CorrespondingEmail}}

%\vfill

% -- Top logo bar
\begin{center}
\includegraphics[height=8cm]{asclab-logo.png}
\end{center}

% -- Date and sponsor
{\large \ReportDate}\\[1cm]
\vspace{1 cm}

\begin{minipage}{0.85\linewidth}
    \centering\small\color{RPIdark}
    \textit{\SponsorText}
\end{minipage}

\vspace{1 cm}

\vspace{0.5cm}
{\color{RPIred}\rule{\linewidth}{2pt}}

\vspace{0.3cm}
{\small\color{RPIdark}
    \textbf{Advanced Space Concepts Laboratory (ASCLab)} $\cdot$
    Rensselaer Polytechnic Institute $\cdot$
    \href{https://www.asclabrpi.com}{asclabrpi.com}
}
\end{center}
\restoregeometry
\end{titlepage}

%% ---- ABSTRACT & KEYWORDS ------------------------------------
\thispagestyle{empty}
\begin{center}
{\large\bfseries\color{ASCblue}\ReportTitle}\\[0.3cm]
{\normalsize \AuthorList}
\end{center}

\begin{abstract}
\noindent
Replace this text with your abstract (150–250 words). The abstract should concisely state
the motivation, problem formulation, methodology, key results, and conclusion.
Avoid using references or unexplained acronyms in the abstract.
\end{abstract}

\noindent\textbf{Keywords:} \Keywords

\vspace{0.5cm}
{\color{RPIred}\rule{\linewidth}{0.8pt}}

%% ---- TABLE OF CONTENTS (optional — comment out if unneeded) -
\tableofcontents
\newpage

%% =============================================================
%%  BODY — edit from here
%% =============================================================

\section{Introduction}
\label{sec:intro}

Provide background, motivation, and a clear statement of the problem. Summarize the
contributions of this report and outline the structure of the document. Cite prior
work using \texttt{\textbackslash citep\{key\}} or \texttt{\textbackslash citet\{key\}}.

Example citation: The CR3BP is a canonical model for cislunar trajectory design
\citep{szebehely1967theory}.

\section{Theoretical Background}
\label{sec:background}

\subsection{Equations of Motion}
\label{sec:eom}

The non-dimensional equations of motion in the CR3BP are
\begin{equation}
    \ddot{x} - 2\dot{y} = \frac{\partial \Omega^*}{\partial x}, \quad
    \ddot{y} + 2\dot{x} = \frac{\partial \Omega^*}{\partial y}, \quad
    \ddot{z}             = \frac{\partial \Omega^*}{\partial z},
    \label{eq:cr3bp}
\end{equation}
where $\Omega^* = \tfrac{1}{2}(x^2+y^2) + \tfrac{1-\mu}{r_1} + \tfrac{\mu}{r_2}$
is the pseudo-potential and $\mu = m_2/(m_1+m_2)$ is the mass ratio.

\subsection{Another Subsection}

Additional theory here.

\section{Methodology}
\label{sec:method}

Describe your algorithms, estimation framework, sensor models, or optimization approach.
Use theorem environments where appropriate.

\begin{theorem}[Jacobi Integral]
\label{thm:jacobi}
The quantity $C = 2\Omega^* - (\dot{x}^2 + \dot{y}^2 + \dot{z}^2)$ is constant along
any solution of \eqref{eq:cr3bp}.
\end{theorem}

\begin{proof}
Take the dot product of the velocity vector with the equations of motion \ldots
\end{proof}

\section{Results and Discussion}
\label{sec:results}

\subsection{Simulation Setup}

Describe your test environment, initial conditions, and parameters.

\subsection{Key Results}

\begin{figure}[H]
    \centering
    % \includegraphics[width=0.75\linewidth]{figures/your_figure.pdf}
    \caption{Replace with your figure. Use PDF or EPS for vector graphics.}
    \label{fig:result1}
\end{figure}

Results can be tabulated as follows.

\begin{table}[H]
    \centering
    \caption{Summary of simulation parameters.}
    \label{tab:params}
    \begin{tabular}{lcc}
        \toprule
        \textbf{Parameter} & \textbf{Symbol} & \textbf{Value} \\
        \midrule
        Mass ratio         & $\mu$           & 0.01215 \\
        Initial epoch      & $t_0$           & J2000.0 \\
        Integration tol.   & ---             & $10^{-12}$ \\
        \bottomrule
    \end{tabular}
\end{table}

\section{Conclusion}
\label{sec:conclusion}

Summarize the main findings and their significance. State limitations and suggest
directions for future work.

%% ---- ACKNOWLEDGMENTS ----------------------------------------
\section*{Acknowledgments}
\addcontentsline{toc}{section}{Acknowledgments}

The authors thank \ldots\ for helpful discussions. \SponsorText

%% ---- REFERENCES ---------------------------------------------
\bibliographystyle{plainnat}    % or aiaa, apalike, etc.
\bibliography{references}       % <-- your .bib file name (without extension)

%% Example manual reference (remove once you have a .bib file):
\begin{thebibliography}{9}
\bibitem{szebehely1967theory}
V.~Szebehely,
\textit{Theory of Orbits: The Restricted Problem of Three Bodies}.
Academic Press, New York, 1967.
\end{thebibliography}

\end{document}
